%! Solving Polynomial Equations by Factoring — Unit 1, Lesson 1.6 — Guided Notes
\documentclass[10pt]{article}
\usepackage{atda-article}
\usepackage{atda-boxes}
\newcommand{\TallMath}[1]{$\displaystyle #1\rule[-1.4em]{0pt}{3.2em}$}
\setlength{\workrowsep}{6pt}

\begin{document}

\pageheader{Unit 1, Lesson 1.6}{Guided Notes}
% NAMESTRIP: no \namedateperiod here. The student writes their name once, on the
% cover this component is stapled behind. See references/conventions.md.

\vspace{0.15in}

\begin{objectivebox}
\small
By the end of this lesson, I will be able to\dots
\begin{enumerate}[leftmargin=1.5em, itemsep=2pt, topsep=3pt]
  \item state the \textbf{zero product property} and explain why only the product \emph{zero}
        tells you anything about the factors;
  \item \textbf{solve a polynomial equation by factoring} --- standard form first, then factor
        completely, then set each factor equal to zero;
  \item say how \textbf{many} solutions a polynomial equation has and what \textbf{kind} they are
        (real or imaginary), and write a factored equation \textbf{from} its solutions or from the
        $x$-intercepts of its graph;
  \item interpret solutions in context and \textbf{reject} the ones the situation forbids.
\end{enumerate}
\end{objectivebox}

\vspace{0.10in}

\boxguard
\begin{vocabbox}
\small
Fill in each term as we build it together.
\termblanklong{Zero product property}
\termblanklong{Solution (root) of an equation}
\termblanklong{Zero of a function}
\termblanklong{Double root}
\termblanklong{Imaginary solution}
\end{vocabbox}

\vspace{0.10in}

\boxguard[12]
\begin{hookbox}
\small
\textbf{The T-shirt launcher.} At the pep rally a shirt is fired from the top of the bleachers,
$48$ feet above the gym floor. Its height above the floor $t$ seconds later is
\[ h(t) = -16t^2+32t+48 \quad \text{feet.} \]
The crew needs to know \textbf{when the shirt hits the floor} --- that is, when $h(t)=0$.

You already have every tool for this. Factor the right-hand side completely, then use last night's
item 13(b): a product is zero only when one of the factors \emph{is} zero. Which $t$ does the story
allow?

\writelines{2}
\end{hookbox}

\vspace{0.10in}

\boxguard[24]
\begin{notesbox}{1. The Zero Product Property --- Why Factoring Solves Anything}
\small
Warm-up item 3 is the whole idea. If $ab=12$ you learn nothing: $a$ could be $1$, $2$, $3$, $12$,
or a fraction. But if $ab=0$, one of the two \emph{has} to be zero --- there is no other way to
make a product vanish.

\vspace{2pt}
\textbf{The zero product property.} If $ab = 0$, then $a=0$ or $b=0$ (or both). It holds for any
number of factors, and \blank{2.2cm} is the only number it works for.

\vspace{2pt}
\textbf{So an equation that is already factored is already almost solved.} Set each factor equal to
zero and solve the small equations that fall out.

\renewcommand{\arraystretch}{1.9}
\begin{tabularx}{\linewidth}{@{}p{3.4cm} p{4.2cm} X@{}}
\rowcolor{mist}
\textbf{Equation} & \textbf{Each factor $= 0$} & \textbf{Solutions} \\
\hline
$(x-4)(x+7)=0$  & \blank{3.6cm} & \blank{2.6cm} \\
$3x(x-5)=0$     & \blank{3.6cm} & \blank{2.6cm} \\
$(2x-1)(x+6)=0$ & \blank{3.6cm} & \blank{2.6cm} \\
$(x-3)^2=0$     & \blank{3.6cm} & \blank{2.6cm} \\
\end{tabularx}

\vspace{6pt}
\textbf{Two things to notice.} In row 2 the constant factor $3$ is never zero, so it contributes no
solution --- but the factor $x$ does, and $x=0$ is a real solution you must keep. In row 4 both
factors are the same, so the two solutions collapse into one: $x=3$ is a \textbf{double root}.
\end{notesbox}

\vspace{0.10in}

\boxguard[24]
\begin{notesbox}{2. Solving by Factoring --- Standard Form First}
\small
The property only works against \textbf{zero}. So before you factor anything, get every term onto
one side:
\begin{enumerate}[leftmargin=1.4em, itemsep=2pt, topsep=3pt]
  \item \textbf{Set the equation equal to zero} (standard form).
  \item \textbf{Factor completely} --- the whole order from Lessons 1.3--1.5, GCF first.
  \item \textbf{Set each factor equal to zero} and solve.
  \item \textbf{Check} each solution in the original equation.
\end{enumerate}

\renewcommand{\arraystretch}{1.9}
\begin{tabularx}{\linewidth}{@{}p{3.2cm} p{4.4cm} X@{}}
\rowcolor{mist}
\textbf{Equation} & \textbf{Standard form, factored} & \textbf{Solutions} \\
\hline
$x^2-7x+12=0$  & \blank{3.8cm} & \blank{2.6cm} \\
$x^2+5x=24$    & \blank{3.8cm} & \blank{2.6cm} \\
$x^2-25=0$     & \blank{3.8cm} & \blank{2.6cm} \\
$2x^2=8x$      & \blank{3.8cm} & \blank{2.6cm} \\
\end{tabularx}

\vspace{6pt}
\textbf{Row 4 is the trap of the day.} It is tempting to divide both sides of $2x^2=8x$ by $x$ and
write $2x=8$, so $x=4$. That answer is true --- and it is not the whole answer, because dividing by
$x$ quietly assumes $x \ne 0$, and $x=0$ \emph{is} a solution here:
\begin{work}
  2x^2 &= 8x \\
  2x^2-8x &= 0 \\
  2x(x-4) &= 0 \\
  x &= 0 \quad \text{or} \quad x = 4 \\
  \text{check } x=0: \; 2(0)^2 &= 8(0) \quad\checkmark \qquad \text{check } x=4: \; 2(16) = 8(4)
      \quad\checkmark
\end{work}
\textbf{Never divide both sides by a variable.} Move it to the other side and let it be a factor.
\end{notesbox}

\vspace{0.10in}

\boxguard[24]
\begin{notesbox}{3. Degree Three and Higher --- How Many, and What Kind}
\small
Nothing changes: standard form, factor completely, set each factor to zero. What is new is that you
can say \emph{in advance} how many solutions to expect --- the \blank{2.0cm} of the polynomial is
the number of solutions, counting a double root twice.
\begin{work}
  x^3-9x &= 0 \\
  x\left(x^2-9\right) &= 0 \\
  x(x+3)(x-3) &= 0 \\
  x &= 0, \quad x = -3, \quad x = 3
\end{work}
Degree $3$, three factors, three solutions --- and the GCF step is what produced the solution $x=0$.
\begin{work}
  x^3+2x^2-15x &= 0 \\
  x\left(x^2+2x-15\right) &= 0 \\
  x(x+5)(x-3) &= 0 \\
  x &= 0, \quad x = -5, \quad x = 3
\end{work}

\vspace{2pt}
\textbf{But not every solution is a real number.} Lesson 1.5 proved that a \emph{sum} of squares
does not factor over the integers. Here is what that costs you:
\begin{work}
  x^3+4x &= 0 \\
  x\left(x^2+4\right) &= 0 \\
  x &= 0 \quad \text{or} \quad x^2+4 = 0 \\
  x^2 &= -4 \quad \text{--- no real number squares to a negative}
\end{work}
So $x=0$ is the only \textbf{real} solution. The equation still has degree $3$, so it still has
three solutions; the other two are \blank{2.4cm} --- numbers outside the real number line, which
you will meet by name in Algebra 2. For now, \textbf{count them and name their type}: one real, two
imaginary.
\end{notesbox}

\vspace{0.10in}

\boxguard[26]
\begin{notesbox}{4. Running the Method Backwards --- From Solutions to an Equation}
\small
Every lesson in this unit has had a backwards direction, and this one is no different. If you know
what the solutions are, you can \textbf{build} the equation: a solution $x=r$ comes from the factor
$(x-r)$, so put one factor in for each solution and set the product equal to zero.

\renewcommand{\arraystretch}{1.9}
\begin{tabularx}{\linewidth}{@{}p{3.4cm} p{4.4cm} X@{}}
\rowcolor{mist}
\textbf{Solutions wanted} & \textbf{Factored equation} & \textbf{Degree} \\
\hline
$5$ and $-3$           & \blank{3.8cm} & \blank{1.8cm} \\
$0$ and $7$            & \blank{3.8cm} & \blank{1.8cm} \\
$-2$, $1$, and $4$     & \blank{3.8cm} & \blank{1.8cm} \\
$0$, $-3$, and $3$     & \blank{3.8cm} & \blank{1.8cm} \\
\end{tabularx}

\vspace{6pt}
\textbf{A graph hands you the same information.} Where a graph crosses the $x$-axis, the height is
zero --- so an $x$-intercept of the graph is a \textbf{zero of the function} and a solution of the
equation. Read the three $x$-intercepts below and write a factored equation with exactly those
solutions.

\begin{center}
\begin{tikzpicture}
\begin{axis}[width=8.2cm, height=4.6cm,
    axis lines=middle, xlabel=$x$, ylabel=$y$,
    xlabel style={font=\scriptsize, at={(ticklabel* cs:1.0)}, anchor=west},
    ylabel style={font=\scriptsize, at={(ticklabel* cs:1.0)}, anchor=south},
    xmin=-3.4, xmax=4.4, ymin=-11, ymax=11,
    xtick={-3,-2,-1,1,2,3,4}, ytick={-8,-4,4,8},
    tick label style={font=\scriptsize},
    grid=both, grid style={line width=0.2pt, draw=linegray},
    samples=101, domain=-2.6:3.5]
  \addplot[royal, very thick, smooth] {x^3 - x^2 - 6*x};
\end{axis}
\end{tikzpicture}
\end{center}

\writelines{2}
\end{notesbox}

\vspace{0.10in}

\boxguard[20]
\begin{practicebox}
\small
\textbf{The second launcher.} A smaller launcher is fired from a $32$-foot platform, so the shirt's
height above the floor is $h(t) = -16t^2+16t+32$ feet.

\begin{enumerate}[leftmargin=1.6em, itemsep=6pt, topsep=4pt]
  \item Factor $h(t)$ completely. (GCF first --- and $-16$ is part of it.)
\begin{work}
  h(t) &= -16t^2+16t+32 \\
       &= -16\left(t^2-t-2\right) \\
       &= -16(t-2)(t+1)
\end{work}

  \item Solve $h(t)=0$. Give both solutions, say which one the story rejects, and why.
\begin{work}
  -16(t-2)(t+1) &= 0 \\
  t-2 &= 0 \quad\Rightarrow\quad t = 2 \\
  t+1 &= 0 \quad\Rightarrow\quad t = -1 \; \text{(rejected --- time cannot be negative)}
\end{work}

% BOXGUARD: \boxguard is inert inside a breakable tcolorbox, and without this
% break the box splits after item 3's prompt, stranding two write-lines alone on
% the last page. Breaking here sends item 3 over whole. Mirrored in notes_key.
\tcbbreak
  \item On its way up and again on its way down, the shirt passes $32$ feet. Find \textbf{both}
        times --- set $h(t)=32$ and put the equation in standard form first --- then say what each
        solution means at the rally.
\begin{work}
  -16t^2+16t+32 &= 32 \\
  -16t^2+16t &= 0 \\
  -16t(t-1) &= 0 \\
  t &= 0 \quad \text{or} \quad t = 1
\end{work}
\writelines{2}
\end{enumerate}
\end{practicebox}

\end{document}
